% This is a template for your written document.
%
% To compile using latexmk on the command line, run the following: 
% latexmk -pdf main.tex

\documentclass[12pt]{article}
\usepackage{setspace}
\singlespace
\usepackage[left=1in,right=1in,top=1in,bottom=1in]{geometry}
\usepackage{graphicx}

\title{\textbf{Project Topic}}
\author{Your Name}

\begin{document}

\maketitle

\section{Introduction}
The introduction will broadly introduce your topic, motivate it, and describe what will be done.

Traffic management remains a crucial issue in large cities, especially rapidly developing cities without extensive urban planning.
Poor traffic management results in congestion, unnecessary emissions, and accidents that impact the safety of the city's residents.
In order to manage the flow of traffic within a city, many cities turn to data-driven approaches that inform urban planning.

"Smart cities" refer to cities that utilize data and computational techniques to improve the functioning of a city. 
For example, a smart city may use location services on cell phones to determine the best times and places to collect waste within a city. 
Intelligent traffic systems, a crucial component of smart cities, aggregate and analyze data to manage traffic within a city. 
Intelligent traffic systems may make use of traffic data, such as data
gathered from Uber or from sensors placed around the city, to model traffic in real time. 
These traffic models can then be used to create software that responds to current traffic flows and decreases congestion, 
such as software like Google Maps that re-routes traffic to less congested areas. Urban planners, public sector transportation 
employees, politicians, and everyday people rely on accurate traffic models to make informed decisions about new transit routes
or traffic policy.

Traffic accident analysis is an important subsection of traffic management. Motor vehicle crashes are the top cause of death
in the United States among people ages 1-44, with traffic fatalities causing about 44,000 deaths in 2022. ~\cite{cdc}
Traffic accident modeling utilizes traffic accident data to determine when and where traffic accidents will occur. This data can 
be used by city planners to create policy geared towards decreasing traffic accidents. For example, [throw in an example of how policymakers use traffic
data to make decisions about traffic, like using smartphone or something.]
Intelligent traffic systems additionally use traffic models to reroute traffic away from accident-prone areas. 

The majority of traffic accident prediction studies focus on determining merely whether or not an accident will happen, but
leave out additional analysis crucial to urban planners or policymakers that help these people to make informed traffic planning
decisions. Traffic accident models which solely predict accidents do not explain how or what factors led into the algorithm
deciding the result, which is both uninformative and seen as unreliable by policymakers who don't understand how the result came to be. 
Thus, many traffic studies rely on explanatory models to analyze the underlying factors that contribute to an accident or the
severity of the accident. Additionally, the majority of traffic accident models focus on large cities where data is more
readily available, rather than small cities where different factors may contribute to accident occurance and severity.

This study aims to provide better insight into traffic accidents in small American towns, including factors that can influence 
the likelihood and severity of a traffic accident. Using XGBoost analysis, 
this study will analyze a dataset of roughly 10,000 traffic accidents in Wooster, Ohio from 2021-2026 to model the 
likelihood of accidents, the severity of the accident, and determine which factors most contributed to causing a traffic accident. 
The results from this study will provide insight into what factors cause traffic accidents, where these accidents, providing
information to traffic officials on strategic decisions that can be made to enhance public safety. 


\section{Background}
A long history of data-driven traffic models that have been used to make planning decisions.
Other approaches involve plotting a simple database of
accident data onto a map to give the data an easily visualized, geospatial component. This approach manages historical accident data
but does not model or predict future accidents.
The field has relied on statistical regression analysis as a predictive model in the past, 
such as historical averages or ARIMA. Historical averages take the average of all data over a certain time period and use that
value as a prediction of what will happen in the future. ARIMA, or Autoregressive Integrated Moving Average, provides a more complex
analysis by integrating three components: "Autoregressive," which computes the current volume of traffic based on previous traffic counts
(ex. if traffic is high at 10 am, it will likely still be high at 10:10 am),
"Integrated," which handles broad trends over long periods of time, and "Moving Average," which allows the model to correct itself
based on its past errors. These methods are computationally simple, but their simplicity
limits their ability to handle non-linear relationships or the incorporation of additional variables, such as weather. 
ARIMA, specifically, handles short-term time-based predictions well, but struggles with real-time traffic modeling, managing
the impact of traffic flows from multiple locations, or handling data with a large variety of variables. 
The accuracy of statistical analysis techniques often cannot measure to the accuracy of
machine learning models and neural networks, which have recently dominanted the traffic modeling field.

Recent approaches using machine learning and neural networks produce models 
of traffic flows, jams, and accidents have shown a marked increase in accuracy of traffic modeling.
A large body of literature showcase the various ways machine learning has been used to model traffic flows.
One study in Thailand models traffic flows and car sales to track emissions across the city, while another study models traffic
to create a signal traffic control system in Ali Mendjeli, Algeria. Others model traffic jams and others focus on accident prediction, but
most traffic studies share similar methods, machine learning techniques, and limitations in the creation of their traffic models. 

Machine learning methods perform well for classification problems, that is, creating a model that will say whether an accident
will occur or not. A plurality of studies have observed the superior performance of decision tree networks such as random forest and XGBoost over neural 
networks such as LSTM or CNNs. Tabular data and often not real-time data.

Neural networks perform well in 


Studies share limitations and concerns. In any smart city project, privacy becomes an issue as personal travel data is being used to
create traffic flows. Data privacy concerns contribute to the already herculean difficulty in acquiring the necessary data to create traffic models.
Even when data is available, the data may not accurately represent the actual traffic patterns in the city. For instance, many studies use data acquired from Uber or Waze
to create neural network-based models, but acknowledge that training a model based on only one transportation platform may be biased since
it would not represent the traffic of all individuals within the city. Finally, data availibility varies between cities and regions, which
results in more data-rich areas receiving more researching resources.
Urbanized cities in North America have a larger abundance of data available to train machine learning models or neural networks, such as Uber 
data, sensors, cameras, and satellite data. Developing cities face greater challenges in acquiring the data required to model traffic flows,
and research projects often require collaboration with local governments in order to place sensors or otherwise acquire traffic data that can be used for analysis.



\section{Related Works}

\section{Methodology}
Dijkstra's algorithm was chosen for single-source shortest path computation on graphs with non-negative edge weights, following the implementation presented in CLRS~\cite{clrsAlgorithms}.

\section{Results}
An exemplary figure is shown in Figure~\ref{fig:dijkstra}. Any time a figure is used, it must have a capation and be directly referred to and discussed in the text.

%\begin{figure}[b] puts the image to the (closest) bottom of page
%\begin{figure}[t] puts the image to the (closest) top of page
%\begin{figure}[h] allows the image to float, putting it roughly wherever it was placed in the text.
\begin{figure}[b] 
\begin{center}
  \includegraphics[width=5cm]{imgs/dijkstra.png}
  \caption{Dijkstra's algorithm does not work on graphs with negative edge weights, as evidenced here.}
  \label{fig:dijkstra}
 \end{center}
\end{figure}



\section{Conclusion}

\bibliographystyle{acm}
\bibliography{bibliography.bib}

\end{document}
